\section{The exact weak-coupling marginal estimate}\label{app:collective}

Let $\Phi_\eta=\Phi_{\tau,n}$. Write $\mathcal C_j(Y)=i[\SW_{0j},Y]$ on the fresh ancilla and block, and let $\mathcal E(X)=\tau\otimes X$, $\mathcal T(Y)=\Tr_0Y$. Then
\begin{equation}
 \Phi_\eta=\mathcal T e^{\eta\mathcal C_n}\cdots e^{\eta\mathcal C_1}\mathcal E.
\end{equation}
On Hermitian operators, each $e^{\eta\mathcal C_j}$ is a trace-norm isometry and each $\mathcal C_j$ has induced norm at most two. Appending $\tau$ is an isometry and partial trace is a contraction on Hermitian inputs. Differentiating the product three times, the sum of multinomial coefficients is $n^3$, so
\begin{equation}
 \norm{\Phi_\eta^{(3)}(X)}_1\le8n^3\norm X_1.
\end{equation}
Taylor's integral remainder therefore gives
\begin{equation}
 \norm{\Phi_\eta(X)-X-\eta^2\mathcal L(X)}_1
 \le\frac43n^3|\eta|^3\norm X_1.
 \label{eq:remainder}
\end{equation}
The first derivative vanishes because the ancilla is maximally mixed. Using $\SW_{0j}=\tfrac12(I+\sum_\alpha\sigma_\alpha^{(0)}\sigma_\alpha^{(j)})$ and $\Tr(\tau\sigma_\alpha\sigma_\beta)=\delta_{\alpha\beta}$, the second-order term is
\begin{equation}
 \mathcal L(X)=-\frac12\sum_{\alpha=1}^3[J_\alpha,[J_\alpha,X]].
 \label{eq:collectivegenerator}
\end{equation}
The cross terms are symmetric after the ancilla trace; the ordering of collisions creates no second-order extra term for this $\tau$ initializer.

Under collective rotations, the operator space splits into irreducible tensor ranks $\ell$. The Casimir in~\eqref{eq:collectivegenerator} acts as $-\ell(\ell+1)/2$. Thus $\mathcal L=-\id$ on the full rank-one isotypic subspace, including its multiplicity spaces. Exact rotational covariance of $\Phi_\eta$ preserves that subspace. If $\eta\le3/(8n^3)$,~\eqref{eq:remainder} gives there
\begin{equation}
 \norm{\Phi_\eta(X)}_1\le(1-\eta^2/2)\norm X_1.
 \label{eq:rankonecontract}
\end{equation}

The rank-one Hilbert--Schmidt projection of $p^{\otimes n}$ lies in the symmetric spin-$n/2$ sector and is
\begin{equation}
 \Pone(p^{\otimes n})=\frac6{(n+1)(n+2)}J_z\Pi_{\rm sym}.
 \label{eq:projection}
\end{equation}
To check the coefficient, the numerator of its projection onto $J_z$ is $n/2$, and $\Tr_{\rm sym}J_z^2=n(n+1)(n+2)/12$. Summing the absolute spin eigenvalues gives trace norm $3n/[2(n+1)]$ for even $n$ and $3(n+1)/[2(n+2)]$ for odd $n$, both at most $3/2$.

Every single-qubit Pauli observable is a rank-one tensor. Orthogonality of the other isotypic sectors and covariance let its expectation be calculated from~\eqref{eq:projection} alone. Iterating~\eqref{eq:rankonecontract} bounds each one-site Bloch expectation by $(3/2)e^{-M\eta^2/2}$. The initial block and its evolution are invariant under rotations about $z$, so each one-qubit marginal is diagonal. Its trace distance from $\tau$ is half the magnitude of that expectation. This proves~\cref{lem:collective}.

The same expansion explains, without replacing the finite-parameter theorem, a correlated weak-coupling limit. At fixed $n$, choose $\eta=M^{-\alpha}$ with $1/3<\alpha<1/2$. The total third-order remainder $O_n(M\eta^3)$ vanishes, while $M\eta^2\to\infty$. A channel telescope comparing $\Phi_\eta^M$ with $e^{M\eta^2\mathcal L}$ then gives
the limit in Eq.~\eqref{eq:collectivetwirl}.
The semigroup is the collective rotational heat semigroup, whose infinite-time limit is the rotation twirl; the initial vector is in the irreducible symmetric sector. The extra comparison error between $I+\eta^2\mathcal L$ and $e^{\eta^2\mathcal L}$ is $O_n(\eta^4)$ per step and also vanishes in the telescope. Meanwhile~\eqref{eq:telescope} tends to zero as $O_n(M^{1/2-2\alpha})$. Every one-site marginal of the limit is $\tau$, but its joint trace distance from $\tau^{\otimes n}$ is $1-(n+1)/2^n$. At $n=2$ this is $1/4$, attaining the zero-return limiting obstruction in~\eqref{eq:jointobstruction}. 

